\documentclass[10pt,a4paper]{article}

% Defense-oriented technical guide for the validated PN-GMP/PNNN comparison.
% Numerical tables and plots are read directly from the frozen result CSV files.
% This document does not execute identification, training, pruning, or inference.

\usepackage[T1]{fontenc}
\usepackage[utf8]{inputenc}
\usepackage[spanish,es-noquoting]{babel}
\usepackage[a4paper,margin=1.65cm,headheight=14pt]{geometry}
\usepackage[protrusion=true,expansion=false]{microtype}
\usepackage{amsmath,amssymb,mathtools}
\usepackage{booktabs,tabularx,longtable,array,multirow,makecell}
\usepackage{float}
\usepackage{xcolor}
\usepackage{enumitem}
\usepackage{listings}
\usepackage{tikz}
\usetikzlibrary{arrows.meta,positioning,calc,fit,shapes.geometric}
\usepackage{pgfplots}
\usepackage{pgfplotstable}
\pgfplotsset{compat=1.18}
\usepackage{hyperref}
\usepackage{fancyhdr}

\definecolor{navy}{HTML}{264653}
\definecolor{teal}{HTML}{2A9D8F}
\definecolor{gold}{HTML}{E9C46A}
\definecolor{orange}{HTML}{F4A261}
\definecolor{redsoft}{HTML}{E76F51}
\definecolor{graybox}{HTML}{F1F3F5}
\definecolor{bluebox}{HTML}{EAF4F4}

\hypersetup{
  colorlinks=true,
  linkcolor=navy,
  urlcolor=teal,
  citecolor=navy,
  pdftitle={Guía de defensa: PN-GMP, DOMP y comparación justa con PNNN},
  pdfauthor={Proyecto PNNN}
}
\pagestyle{fancy}
\fancyhf{}
\lhead{PN-GMP y PNNN: guía de defensa}
\rhead{Resultados validados 20260714\_102201}
\cfoot{\thepage}
\setlength{\parindent}{0pt}
\setlength{\parskip}{3.5pt}
\setlist{nosep,leftmargin=1.4em}
\setlength{\emergencystretch}{2em}
\raggedbottom

\newcommand{\RunDir}{../results/full_signal_domp_comparison/20260714_102201}
\pgfplotstableread[col sep=comma]{\RunDir/complexity_flops.csv}\ComplexityTable

% Verified numerical projection of comparison_results.csv.  The source CSV is
% RFC-quoted because one text column contains commas; pgfplotstable does not
% parse those quoted commas reliably, so only the used numeric fields are kept
% here.  The scientific result file remains unchanged.
\newcommand{\Rdef}[3]{\expandafter\def\csname R@#1@#2\endcsname{#3}}
\newcommand{\Rval}[3]{%
  \edef\ResultValue{\csname R@#1@#2\endcsname}%
  \pgfmathprintnumber[fixed,precision=#3]{\ResultValue}}
\Rdef{0}{NumRealParameters}{200}
\Rdef{0}{InternalTrainNMSEdB}{-39.0822317101167}
\Rdef{0}{InternalValidationNMSEdB}{-38.3905172043167}
\Rdef{0}{IdentificationNMSEdB}{-39.0398273701677}
\Rdef{0}{FullSignalNMSEdB}{-37.9778360655445}
\Rdef{1}{NumRealParameters}{200}
\Rdef{1}{InternalTrainNMSEdB}{-39.0822317101175}
\Rdef{1}{InternalValidationNMSEdB}{-38.3905172046381}
\Rdef{1}{IdentificationNMSEdB}{-39.0398273701892}
\Rdef{1}{FullSignalNMSEdB}{-37.9778360655442}
\Rdef{1}{RelativePredictionErrorToComplex}{1.32062424670276e-11}
\Rdef{2}{NumRealParameters}{358}
\Rdef{2}{EffectiveFeatureCount}{179}
\Rdef{2}{InternalTrainNMSEdB}{-39.7334782614176}
\Rdef{2}{InternalValidationNMSEdB}{-38.9055344784472}
\Rdef{2}{IdentificationNMSEdB}{-39.7239821395781}
\Rdef{2}{FullSignalNMSEdB}{-38.5336473098535}
\Rdef{3}{NumRealParameters}{358}
\Rdef{3}{EffectiveFeatureCount}{179}
\Rdef{3}{InternalTrainNMSEdB}{-39.4283004030278}
\Rdef{3}{InternalValidationNMSEdB}{-38.5727743709773}
\Rdef{3}{IdentificationNMSEdB}{-39.384614952974}
\Rdef{3}{FullSignalNMSEdB}{-38.1702651011503}
\Rdef{4}{NumRealParameters}{394}
\Rdef{4}{EffectiveFeatureCount}{197}
\Rdef{4}{InternalTrainNMSEdB}{-39.3621947321484}
\Rdef{4}{InternalValidationNMSEdB}{-38.00618750765}
\Rdef{4}{IdentificationNMSEdB}{-39.2580287899067}
\Rdef{4}{FullSignalNMSEdB}{-37.8104589405615}
\Rdef{5}{NumRealParameters}{200}
\Rdef{5}{EffectiveFeatureCount}{100}
\Rdef{5}{InternalTrainNMSEdB}{-39.0864733714523}
\Rdef{5}{InternalValidationNMSEdB}{-38.2995281181917}
\Rdef{5}{IdentificationNMSEdB}{-39.0251889044127}
\Rdef{5}{FullSignalNMSEdB}{-38.0578594656762}
\Rdef{6}{NumRealParameters}{350}
\Rdef{6}{InternalTrainNMSEdB}{-34.2382052677646}
\Rdef{6}{InternalValidationNMSEdB}{-33.5119879920964}
\Rdef{6}{IdentificationNMSEdB}{-34.281119051028}
\Rdef{6}{FullSignalNMSEdB}{-33.6332342790055}
\Rdef{7}{NumRealParameters}{1046}
\Rdef{7}{InternalTrainNMSEdB}{-37.5542109467709}
\Rdef{7}{InternalValidationNMSEdB}{-36.6302846274961}
\Rdef{7}{IdentificationNMSEdB}{-37.4845848646783}
\Rdef{7}{FullSignalNMSEdB}{-36.6310229561146}
\Rdef{8}{NumRealParameters}{358}
\Rdef{8}{InternalTrainNMSEdB}{-36.2817322310163}
\Rdef{8}{InternalValidationNMSEdB}{-35.5243894899996}
\Rdef{8}{IdentificationNMSEdB}{-36.5743240001578}
\Rdef{8}{FullSignalNMSEdB}{-35.8997812260251}
\newcommand{\Cval}[3]{%
  \pgfplotstablegetelem{#1}{#2}\of{\ComplexityTable}%
  \pgfmathprintnumber[fixed,precision=#3]{\pgfplotsretval}}

% Coordinates used in plots are extracted from the validated CSV at compile time.
\xdef\PNParams{358}
\xdef\PNFull{-38.5336473098535}
\xdef\GMPMParams{358}
\xdef\GMPMFull{-38.1702651011503}
\xdef\SparseParams{358}
\xdef\SparseFull{-35.8997812260251}
\pgfplotstablegetelem{2}{FLOPsPerSample}\of{\ComplexityTable}\xdef\PNFlops{\pgfplotsretval}
\pgfplotstablegetelem{3}{FLOPsPerSample}\of{\ComplexityTable}\xdef\GMPMFlops{\pgfplotsretval}
\pgfplotstablegetelem{8}{FLOPsPerSample}\of{\ComplexityTable}\xdef\SparseFlops{\pgfplotsretval}
\pgfplotstablegetelem{8}{DenseMatrixFLOPsPerSample}\of{\ComplexityTable}\xdef\SparseDenseMatrixFlops{\pgfplotsretval}

\newenvironment{keybox}{%
  \begin{center}\begin{minipage}{0.94\linewidth}\color{navy}%
  \hrule\vspace{5pt}\color{black}}{%
  \vspace{5pt}\color{navy}\hrule\color{black}\end{minipage}\end{center}}

\lstset{
  basicstyle=\ttfamily\small,
  frame=single,
  backgroundcolor=\color{graybox},
  columns=fullflexible,
  keepspaces=true,
  showstringspaces=false
}

\title{\textbf{PN-GMP, DOMP y comparación justa con PNNN}\\
\large Manual técnico y guía para la defensa oral}
\author{Proyecto PNNN}
\date{Ejecución validada: \texttt{20260714\_102201}\\14 de julio de 2026}

\begin{document}
\maketitle
\begin{center}
\begin{minipage}{0.86\linewidth}
Este documento explica el estudio final sin regenerar resultados. Las cifras se
leen directamente de los CSV congelados de la ejecución indicada. La convención
local se conserva: $X$ es la entrada y $Y$ la salida del bloque modelado;
\texttt{xy\_forward} no se interpreta por sí solo como PA-forward.
\end{minipage}
\end{center}
\vfill
\begin{keybox}
\textbf{Mensaje central.} La formulación PN-IQ acoplada es una representación
real exactamente equivalente al GMP complejo. La mejora de NMSE aparece al
relajar el acoplamiento I/Q en coordenadas normalizadas en fase. En esta captura,
Independent PN-IQ alcanza el menor NMSE sobre la señal completa; una PNNN N12
podada iguala exactamente sus 358 parámetros activos y necesita menos FLOPs si
la implementación omite los pesos nulos. No existe un ganador universal sin
fijar plataforma y criterio.
\end{keybox}
\vfill
\tableofcontents
\clearpage

\section{Resumen ejecutivo}

El objetivo experimental es comparar modelos de señal compleja con memoria en
un protocolo común y con presupuestos de parámetros explícitos. El punto de
partida es un Generalized Memory Polynomial (GMP): una población de regresores
complejos combina muestras actuales y retardadas con términos de envolvente.
DOMP selecciona un soporte disperso durante la identificación. Después, el
modelo inferido solo evalúa las columnas seleccionadas y sus coeficientes; DOMP
no forma parte de la inferencia.

La normalización de fase usa
\begin{equation}
 r(n)=\frac{x^*(n)}{|x(n)|},\qquad r(n)=1\ \text{si}\ x(n)=0,
 \label{eq:rotation}
\end{equation}
y rota con el mismo $r(n)$ toda la fila del GMP y el objetivo. La formulación
\emph{coupled PN-IQ} transforma el sistema complejo en una matriz real por
bloques, conservando exactamente dos parámetros reales por coeficiente complejo.
El error relativo de predicción frente al GMP es
$\Rval{1}{RelativePredictionErrorToComplex}{12}$: equivalencia numérica.

La variante \emph{independent PN-IQ} conserva las features reales normalizadas,
pero ajusta I y Q con coeficientes independientes. Ya no impone la estructura de
una multiplicación compleja y, por tanto, es una clase más general. Los controles
sin PN, con GMP parameter-matched y con PN-IQ reducido separan el efecto de las
coordenadas, el acoplamiento y el número de parámetros.

La captura tiene 491,520 muestras. La identificación contiene 19,661
(aproximadamente 4\%) y está incluida en la evaluación sobre la señal completa.
Dentro de identificación, 16,711 muestras de entrenamiento interno y 2,950 de
validación interna sirven exclusivamente para seleccionar $\lambda$, épocas y
duración del fine-tuning. Después, todos los soportes, coeficientes y redes
finales se reajustan desde cero usando las 19,661 muestras. La señal completa no
interviene en selección.

Independent PN-IQ obtiene \Rval{2}{FullSignalNMSEdB}{3} dB con 358 parámetros;
el GMP parameter-matched, \Rval{3}{FullSignalNMSEdB}{3} dB con 358; y PNNN N12
sparse, \Rval{8}{FullSignalNMSEdB}{3} dB con 358. La primera mejora al GMP matched
en 0.363 dB y a la PNNN sparse en 2.634 dB. La red sparse, a su vez, mejora
2.267 dB a H4 dense, pero empeora 0.731 dB respecto a N12 dense. PN-IQ usa
\Cval{2}{FLOPsPerSample}{0} FLOPs/muestra; N12 sparse usa
\Cval{8}{FLOPsPerSample}{0} cuando se omiten pesos cero, o
\Cval{8}{DenseMatrixFLOPsPerSample}{0} si se ejecutan las matrices completas,
además de 12 evaluaciones sigmoidales y otras operaciones adicionales. La evaluación completa
incluye las muestras de identificación y no es un test independiente de
generalización.

\begin{figure}[ht]
\centering
\begin{tikzpicture}[node distance=7mm and 8mm,>=Latex,
 box/.style={draw=navy,rounded corners,fill=bluebox,align=center,minimum height=8mm,minimum width=25mm,font=\small},
 data/.style={draw=teal,rounded corners,fill=teal!10,align=center,minimum height=8mm,font=\small}]
\node[data] (cap) {Captura $X/Y$\\DC eliminado};
\node[box,right=of cap] (split) {Identificación 4\%\\selección interna};
\node[box,right=of split] (features) {Features\\GMP y PN-IQ};
\node[box,above right=5mm and 9mm of features] (linear) {DOMP + $\lambda$\\modelos lineales};
\node[box,below right=5mm and 9mm of features] (nn) {PNNN dense\\y pruning N12};
\node[box,right=32mm of features] (eval) {Señal completa\\NMSE, parámetros, FLOPs};
\draw[->,thick] (cap)--(split);
\draw[->,thick] (split)--(features);
\draw[->,thick] (features)--(linear);
\draw[->,thick] (features)--(nn);
\draw[->,thick] (linear)--(eval);
\draw[->,thick] (nn)--(eval);
\end{tikzpicture}
\caption{Flujo validado. La señal completa contiene la identificación y no selecciona hiperparámetros.}
\label{fig:flow}
\end{figure}

\section{Contexto de DPD y GMP}

La predistorsión digital busca compensar no linealidad y memoria antes de la
cadena física, pero este repositorio formula cada identificación con la
convención local entrada $X$/salida $Y$ del bloque concreto. Un PA con memoria no
depende solo de $x(n)$: depende también de retardos y envolventes recientes. Un
GMP representa esa dependencia mediante ramas alineadas y cruzadas. En forma
matricial,
\begin{equation}
 \mathbf y=\mathbf U\mathbf h+\mathbf e,
\end{equation}
donde cada columna de $\mathbf U$ es un regresor complejo y cada elemento de
$\mathbf h$ es su coeficiente complejo. Las ramas A/B/C combinan un carrier
retardado con una potencia de envolvente actual, adelantada o retardada; los
builders del proyecto conservan además los auxiliares definidos por la población.

Seleccionar 100 regresores complejos significa mantener 100 columnas complejas,
no 100 escalares reales. Si $h_j=\alpha_j+j\beta_j$, cada coeficiente aporta dos
componentes reales. Por ello,
\begin{equation}
 100\ \text{coeficientes complejos}=200\ \text{parámetros reales}.
\end{equation}
Esta contabilidad es la que hace posible una comparación parameter-matched con
modelos reales y redes.

\section{Phase normalization}

Para cada fila evaluada se calcula \eqref{eq:rotation} y se aplica
\begin{equation}
 \mathbf U_{\mathrm{PN}}(n,:)=r(n)\mathbf U(n,:),\qquad
 y_{\mathrm{PN}}(n)=r(n)y(n).
 \label{eq:pnrow}
\end{equation}
El caso $x(n)=0$ usa $r(n)=1$ y evita división por cero. La implementación valida
índices y finitud y devuelve una rotación columna.

La misma rotación debe multiplicar \emph{toda la fila}. No se construye primero
$z(k)=r(k)x(k)$ para introducirlo en un GMP convencional: un tap $x(n-m)$ quedaría
rotado con $r(n-m)$, cuando el marco deseado es el de la muestra actual $r(n)$.
Para un regresor
\begin{equation}
 u(n)=x(n-m)E(n),\qquad x(k)=|x(k)|e^{j\phi(k)},
\end{equation}
con $E(n)$ real, se obtiene
\begin{equation}
 r(n)u(n)=|x(n-m)|E(n)e^{j[\phi(n-m)-\phi(n)]}.
 \label{eq:relativephase}
\end{equation}
Se elimina la fase absoluta del carrier de referencia, pero se conserva amplitud,
memoria y fase relativa. Si $m=0$, la diferencia de fase es cero y la componente
Q de esa feature es exactamente nula. Solo esas columnas Q estructuralmente
nulas se eliminan en las variantes independientes; ello no convierte el
coeficiente del GMP acoplado en real.

La extensión temporal del constructor compartido usa índices periódicos para
que todas las variantes observen el mismo dominio. Tras predecir en el marco PN,
la salida principal se restaura como
\begin{equation}
 \widehat y(n)=r^*(n)\widehat y_{\mathrm{PN}}(n).
\end{equation}
PN no es necesario para escribir una matriz compleja en forma real. Su papel es
elegir coordenadas relativas a la fase actual antes de relajar el modelo I/Q.

\section{Formulación PN-IQ acoplada}

Sea $u=u_I+ju_Q$ y $h=\alpha+j\beta$. Entonces
\begin{equation}
 uh=(u_I\alpha-u_Q\beta)+j(u_I\beta+u_Q\alpha).
\end{equation}
Las dos vistas equivalentes son
\begin{equation}
 \begin{bmatrix}y_I\\y_Q\end{bmatrix}
 =\begin{bmatrix}u_I&-u_Q\\u_Q&u_I\end{bmatrix}
 \begin{bmatrix}\alpha\\\beta\end{bmatrix},\qquad
 \begin{bmatrix}y_I\\y_Q\end{bmatrix}
 =\begin{bmatrix}\alpha&-\beta\\\beta&\alpha\end{bmatrix}
 \begin{bmatrix}u_I\\u_Q\end{bmatrix}.
 \label{eq:coupled1}
\end{equation}
Para múltiples regresores, con $\mathbf U_{\mathrm{PN}}=\mathbf U_I+j\mathbf U_Q$,
\begin{equation}
 \begin{bmatrix}\mathbf y_I\\\mathbf y_Q\end{bmatrix}
 =\underbrace{\begin{bmatrix}\mathbf U_I&-\mathbf U_Q\\
                 \mathbf U_Q& \mathbf U_I\end{bmatrix}}_{\mathbf U_R}
  \begin{bmatrix}\boldsymbol\alpha\\\boldsymbol\beta\end{bmatrix}.
 \label{eq:coupledmulti}
\end{equation}
Se llama \emph{acoplada} porque los cuatro bloques no son independientes: la
misma $\alpha_j$ y $\beta_j$ gobiernan simultáneamente ambas salidas y aparecen
con signos fijados. Hay dos grados de libertad por regresor, exactamente los de
$h_j$. La rotación de fila es unitaria, de modo que el ajuste y la restauración
producen el mismo GMP salvo redondeo. El resultado validado es
$\Rval{1}{RelativePredictionErrorToComplex}{12}$ de error relativo sobre la
señal completa y el NMSE coincide a la precisión mostrada.

Una implementación \emph{structure-aware} puede evitar operar con las Q
exactamente nulas y compartir envolventes y potencias. Es un ahorro de
operaciones, no una licencia para eliminar $\beta_j$: una feature real aún admite
coeficiente complejo y conserva dos parámetros. El valor de la formulación
acoplada es servir como puente algebraico, prueba de equivalencia y base común
para estudiar restricciones más flexibles.

\begin{figure}[ht]
\centering
\begin{tikzpicture}[>=Latex,font=\small,
 block/.style={draw,rounded corners,minimum width=46mm,minimum height=20mm,align=center},
 note/.style={align=center,text width=43mm}]
\node[block,fill=bluebox] (c) {\textbf{Coupled PN-IQ}\\
$\left[\begin{smallmatrix}U_I&-U_Q\\U_Q&U_I\end{smallmatrix}\right]$
$\left[\begin{smallmatrix}\alpha\\\beta\end{smallmatrix}\right]$};
\node[block,fill=gold!22,right=18mm of c] (i) {\textbf{Independent PN-IQ}\\
$y_I=F a$\\$y_Q=F b$};
\node[note,below=4mm of c] {2 grados/regresor\\equivalente al GMP};
\node[note,below=4mm of i] {coeficientes de salida libres\\clase más general};
\draw[->,thick] ($(c.east)+(0,3mm)$)--node[above]{relajar acoplamiento}($(i.west)+(0,3mm)$);
\draw[->,thick] ($(i.west)+(0,-5mm)$)--node[below]{imponer estructura}($(c.east)+(0,-5mm)$);
\end{tikzpicture}
\caption{La diferencia esencial no es ``complejo frente a real'', sino la estructura impuesta entre I y Q.}
\label{fig:coupled-independent}
\end{figure}

\section{Formulación PN-IQ independiente}

Una transformación real general de una feature I/Q es
\begin{equation}
 \begin{bmatrix}y_I\\y_Q\end{bmatrix}
 =\begin{bmatrix}a&b\\c&d\end{bmatrix}
 \begin{bmatrix}u_I\\u_Q\end{bmatrix}.
 \label{eq:independent2x2}
\end{equation}
Ya no se exige $d=a$ ni $c=-b$ (según convención de signos). En el caso
multifeature se forma una matriz real $\mathbf F$ con las componentes I/Q
retenidas y se ajustan dos salidas:
\begin{equation}
 \mathbf y_I=\mathbf F\mathbf a,\qquad
 \mathbf y_Q=\mathbf F\mathbf b.
\end{equation}
Una feature genérica puede contribuir con cuatro grados de libertad a la pareja
de salidas, frente a dos en una multiplicación compleja. Por eso el modelo puede
reducir NMSE, pero también sobreparametrizar y dejar de ser un GMP equivalente.

\begin{table}[ht]
\centering\small
\caption{Controles que separan coordenadas, acoplamiento y presupuesto.}
\begin{tabularx}{\linewidth}{@{}lXcccc@{}}
\toprule
Modelo & Entrada y salida & PN & Acoplamiento & Parámetros & Interpretación\\
\midrule
Coupled PN-IQ & $U_{\rm PN}\to(y_I,y_Q)$ & sí & complejo & 200 & GMP exacto\\
Independent PN-IQ full & $F_{\rm PN}\to y_I,y_Q$ & sí & libre & 358 & extensión principal\\
Independent I/Q no PN & $[\Re U,\Im U]\to y_I,y_Q$ & no & libre & 394 & control de coordenadas\\
Reduced independent PN-IQ & 100 features compartidas & sí & libre & 200 & control de presupuesto\\
Complex GMP matched & columnas GMP complejas & no & complejo & 358 & control parameter-matched\\
\bottomrule
\end{tabularx}
\end{table}

En el soporte DOMP-100, las 200 columnas I/Q brutas pierden 21 columnas Q
exactamente nulas; quedan 179 features y dos coeficientes de salida por feature:
358 parámetros. El control sin PN conserva 197 features, es decir, 394
parámetros. La versión reducida selecciona con DOMP 100 features reales
compartidas y les asigna dos salidas: 200 parámetros. El GMP matched usa 179
regresores complejos, 358 parámetros: coincide exactamente con el objetivo
congelado del proceso de selección.

\section{Por qué PN importa en el modelo independiente}

Una multiplicación compleja es equivariante a rotaciones. En el acoplado,
\begin{center}
rotación $\rightarrow$ multiplicación compleja $\rightarrow$ rotación inversa
$=$ mismo GMP.
\end{center}
Por eso PN no cambia su clase ni su NMSE. Una matriz real general como
\eqref{eq:independent2x2} no conmuta necesariamente con todas las rotaciones. Al
romper el acoplamiento,
\begin{center}
rotación $\rightarrow$ transformación I/Q general $\rightarrow$ rotación inversa
$=$ modelo dependiente del marco PN.
\end{center}
El control Independent I/Q without PN tiene la misma libertad conceptual sin
elegir ese marco. Su full-signal NMSE es \Rval{4}{FullSignalNMSEdB}{3} dB,
frente a \Rval{2}{FullSignalNMSEdB}{3} dB con PN: en esta captura, PN aporta
0.723 dB dentro de la familia independiente. Esto aísla un efecto de coordenadas, no prueba una
ventaja universal.

\section{DOMP}

La población GMP es mayor que el soporte inferido. La selección dispersa busca
columnas informativas sin evaluar todas en el modelo final. OMP selecciona por
correlación y actualiza el residuo; DOMP añade una ortogonalización explícita de
los candidatos respecto al subespacio ya elegido, reduciendo selecciones
redundantes en diccionarios correlacionados. La versión final usa solo DOMP; OMP
se excluyó para mantener una única regla validada y comparable.

\begin{lstlisting}[caption={Pseudocódigo resumido de DOMP.}]
Z <- normalize_columns(X); residual <- y; support <- []
repeat until maxSupport or stopping criterion:
    score <- FrobeniusNorm(Z' * residual, over outputs)
    j <- largest eligible score
    support <- support union {j}
    q <- normalized orthogonal candidate Z(:,j)
    Z <- Z - q * (q' * Z)              % candidate re-orthogonalization
    coefficients <- min_norm_LS(X(:,support), y)
    residual <- y - X(:,support) * coefficients
end
\end{lstlisting}

\texttt{selectDOMPSupport.m} normaliza columnas, excluye normas por debajo de
tolerancia, impide índices repetidos, acepta matrices y targets complejos o
multisalida, puntúa con norma de Frobenius, y ajusta sobre las columnas
\emph{originales}. Usa QR si hay rango completo y \texttt{lsqminnorm} en caso
contrario. Guarda en \texttt{history} índice, correlación, norma del residuo,
condición, tiempo y ortogonalidad. Las tolerancias de columna, correlación y
residuo, junto con el máximo de soporte, definen la parada. Dada la misma matriz,
target y configuración, el proceso es determinista.

DOMP actúa solo dentro de identificación para elegir soporte (y en el flujo
parameter-matched o reducido correspondiente). Tras la selección interna, el
soporte final se regenera con las 19,661 muestras de identificación. El soporte
DOMP-100 congelado se reutiliza en el coupled. No se suma su coste a
FLOPs/muestra de inferencia. El algoritmo no se presenta como
aportación propia; la referencia de proyecto es Becerra et al., ``A Doubly
Orthogonal Matching Pursuit Algorithm for Sparse Predistortion of Power
Amplifiers'' \cite{becerra_domp}. Los artefactos disponibles no fijan más datos
bibliográficos, por lo que no se inventan año ni revista.

\section{Diseño de la comparación justa}

La medida configurada es \texttt{experiment20260429T134032\_xy}, con 491,520
muestras y eliminación de DC. El mapping guardado es \texttt{xy\_forward}, que
solo identifica la convención local $X/Y$. El constructor aplica la misma
extensión temporal periódica a los modelos con memoria.

Con semilla 1004, el selector original elige aproximadamente el 4\% para
identificación. Con semilla 42, un split estratificado por amplitud de $|Y|$
(10 bins) divide ese dominio en entrenamiento y validación internos:
\begin{equation*}
|\mathcal I_{\rm train}|=16,711,\quad |\mathcal I_{\rm val}|=2,950,\quad
|\mathcal I|=19,661,\quad |\mathcal F|=491,520.
\end{equation*}
Los dos subconjuntos internos no se intersectan y su unión es identificación.
El dominio completo $\mathcal F$ contiene todos los índices de identificación.
Por tanto, la evaluación sobre la señal completa no es un hold-out disjunto.

\begin{figure}[ht]
\centering
\begin{tikzpicture}[font=\small]
\def\W{13.5}
\fill[navy!28] (0,0) rectangle (\W,0.65);
\fill[teal!65] (0,0) rectangle (0.54,0.65);
\draw (0,0) rectangle (\W,0.65);
\node at (7.0,0.325) {Señal completa: 491,520 (incluye identificación)};
\node[align=center] (poollabel) at (1.75,1.15) {Identificación: 19,661\\(4.00004\%)};
\draw[->] (poollabel.south west)--(0.27,0.67);
\node[anchor=east] at (0,-0.9) {Selección interna:};
\fill[teal!55] (0.35,-1.25) rectangle (5.45,-0.65);
\fill[gold!75] (5.45,-1.25) rectangle (6.35,-0.65);
\draw (0.35,-1.25) rectangle (6.35,-0.65);
\node[align=center] at (2.90,-0.95) {Entrenamiento: 16,711};
\node[align=center,font=\scriptsize] at (5.90,-0.95) {Validación\\2,950};
\end{tikzpicture}
\caption{Identificación incluida en la señal completa. La barra inferior amplía su división interna de selección.}
\label{fig:split}
\end{figure}

Para modelos lineales, el subconjunto interno de entrenamiento selecciona
soporte y se compara en validación interna la rejilla
\[
\lambda\in\{0,10^{-6},10^{-5},10^{-4},10^{-3},10^{-2}\}.
\]
La ejecución validada elige $\lambda=0$ en todas las variantes. Congeladas las
decisiones, DOMP vuelve a generar cada soporte final y los coeficientes se
ajustan con identificación completa.

Para PNNN, el split interno selecciona 3,083 épocas para H4, 3,150 para N12 y
392 épocas de fine-tuning. Cada red final se reinicia con la misma semilla,
calcula su normalización sobre identificación completa y ejecuta exactamente el
número de épocas seleccionado sin early stopping. N12 sparse se poda y se
afina también sobre identificación completa. La señal completa solo se predice
cuando los modelos están congelados.

La comparación anterior no era justa si mezclaba identificación al 4\%, split
histórico 70/15/15, una red de unos 1200 parámetros y conjuntos de evaluación
distintos. El protocolo actual comparte identificación, evaluación completa,
función NMSE y cuenta parámetros reales activos. El objetivo sparse se calcula
dinámicamente para igualar los 358 parámetros de Independent PN-IQ. Aun así,
persisten diferencias de optimización: los modelos lineales resuelven ajustes
convexos y la red sparse parte de una arquitectura densa sobredimensionada.

El experimento disjunto anterior se conserva en
\path{results/fair_domp_comparison/20260714_013938} como control exploratorio
histórico, pero no es el resultado principal.

\section{PNNN dense y sparse}

La PNNN recibe $D=84$ features phase-normalized y produce dos salidas reales que
se recombinan y se restauran a la fase original. Tiene una capa oculta sigmoide y una
salida lineal de dos componentes. Para $H$ neuronas,
\begin{equation}
 P(H)=DH+H+2H+2=H(D+3)+2.
\end{equation}
Así, H4 contiene 350 parámetros y H12, 1046. Ambas usan seed 42. H4 dense sirve
como control denso cercano a 358 parámetros; N12 dense mide la capacidad de la
arquitectura ancha; N12 sparse empieza con esa capacidad y se poda hasta 358
parámetros activos.

La selección interna fija 3,083 épocas para H4 y 3,150 para N12. El ajuste final
parte de cero y usa las 19,661 muestras en cada una de esas épocas. En N12
sparse, el pruning global por magnitud elimina 688 de 1032 pesos (66.667\%),
protege los 14 biases, conserva 344 pesos y totaliza 358 parámetros. Los pesos
podados quedan congelados y el modelo se afina durante 392 épocas fijas usando
identificación completa. Proteger biases evita que el presupuesto desactive
offsets accidentalmente; podarlos tendría otra semántica y contabilidad.

Una red ancha sparse puede conservar rutas útiles distribuidas que una red H4 no
puede representar, aunque no hay garantía: depende de inicialización,
optimización y patrón. N12 sparse logra \Rval{8}{FullSignalNMSEdB}{3} dB, mejora H4 en
2.267 dB, pero N12 dense alcanza \Rval{7}{FullSignalNMSEdB}{3} dB y es 0.731 dB mejor
que la podada.

\begin{figure}[ht]
\centering
\begin{tikzpicture}[>=Latex,node distance=12mm,font=\small,
 box/.style={draw=navy,rounded corners,fill=bluebox,minimum width=35mm,minimum height=14mm,align=center}]
\node[box] (dense) {N12 dense\\1046 parámetros\\señal completa \Rval{7}{FullSignalNMSEdB}{3} dB};
\node[box,right=of dense] (prune) {pruning global\\$-688$ pesos\\biases protegidos};
\node[box,right=of prune] (sparse) {N12 sparse\\358 activos\\señal completa \Rval{8}{FullSignalNMSEdB}{3} dB};
\draw[->,thick] (dense)--node[above]{magnitud}(prune);
\draw[->,thick] (prune)--node[above]{freeze + fine-tune}(sparse);
\end{tikzpicture}
\caption{La red sparse conserva la arquitectura H12, no se sustituye por una red estrecha.}
\label{fig:pruning}
\end{figure}

\section{Métrica NMSE}

Todos los modelos usan la misma definición temporal compleja:
\begin{equation}
 \mathrm{NMSE}_{\mathrm{dB}}=10\log_{10}
 \frac{\sum_n|y(n)-\widehat y(n)|^2}{\sum_n|y(n)|^2}.
 \label{eq:nmse}
\end{equation}
Más negativo significa menor energía de error relativa. Una ventaja de 0.3 dB
es pequeña pero medible sobre el mismo dominio; 1 dB es más clara; 3 dB equivale a
aproximadamente dividir por dos la potencia de error. Estas interpretaciones no
sustituyen repetición estadística. No se mezclan NMSE temporal y frecuencial ni
dominios distintos. Full-signal NMSE incluye identificación y no cuantifica por
sí solo generalización independiente.

\section{Parámetros y comparación parameter-matched}

\begin{table}[ht]
\centering\small
\caption{Presupuesto real y soporte de la ejecución validada.}
\begin{tabular}{@{}lrrrr@{}}
\toprule
Modelo & Coef. complejos & Features efectivas & Parámetros reales & Dif. objetivo\\
\midrule
Complex GMP DOMP-100 & 100 & 100 & \Rval{0}{NumRealParameters}{0} & --\\
Coupled PN-IQ & 100 & 200 bloques I/Q & \Rval{1}{NumRealParameters}{0} & --\\
Independent PN-IQ full & -- & \Rval{2}{EffectiveFeatureCount}{0} & \Rval{2}{NumRealParameters}{0} & 0\\
Complex GMP matched & 179 & \Rval{3}{EffectiveFeatureCount}{0} & \Rval{3}{NumRealParameters}{0} & 0\\
Independent I/Q no PN & -- & \Rval{4}{EffectiveFeatureCount}{0} & \Rval{4}{NumRealParameters}{0} & $+36$\\
Reduced independent PN-IQ & -- & \Rval{5}{EffectiveFeatureCount}{0} & \Rval{5}{NumRealParameters}{0} & $-158$\\
PNNN H4 dense & -- & -- & \Rval{6}{NumRealParameters}{0} & $-8$\\
PNNN N12 dense & -- & -- & \Rval{7}{NumRealParameters}{0} & $+688$\\
PNNN N12 sparse & -- & -- & \Rval{8}{NumRealParameters}{0} & 0\\
\bottomrule
\end{tabular}
\end{table}

El objetivo de 358 proviene de las 179 features PN efectivas por dos salidas. El
GMP parameter-matched selecciona exactamente 179 coeficientes complejos, es
decir, 358 parámetros reales. En redes se cuentan pesos y biases activos; N12
sparse tiene 344+14=358. Los tres modelos principales quedan, por tanto,
exactamente igualados en este presupuesto.

\section{FLOPs y complejidad}

Se usa una única cifra principal: \textbf{FLOPs/muestra}, definida como el número
de sumas y multiplicaciones requerido para producir una muestra compleja de
salida a partir de sus muestras de entrada. El alcance es el mismo para todos los
modelos: generación de las features utilizadas, productos por coeficientes o
pesos, acumulaciones, biases, normalización de fase y restauración cuando
corresponden.

La convención asigna 1 FLOP a una suma/resta o multiplicación real, 2 a una suma
compleja y 6 a una multiplicación compleja. Para
\begin{equation*}
(a+jb)(c+jd)=(ac-bd)+j(ad+bc)
\end{equation*}
hay cuatro multiplicaciones y dos sumas reales: 6 FLOPs. Se informan por separado
módulos, raíces, divisiones, evaluaciones de activación y exponenciales porque no se ha acordado una
conversión homogénea a FLOPs.

\begin{table}[ht]
\centering\small
\caption{Comparación principal: mismo presupuesto de 358 parámetros reales.}
\begin{tabularx}{\linewidth}{@{}XrrrX@{}}
\toprule
Modelo & Parámetros & Full-signal NMSE & FLOPs/muestra & Operaciones adicionales/muestra\\
\midrule
Independent PN-IQ & \Rval{2}{NumRealParameters}{0} & \Rval{2}{FullSignalNMSEdB}{3} dB & \Cval{2}{FLOPsPerSample}{0} & 12 módulos/raíces, 2 divisiones\\
Complex GMP matched & \Rval{3}{NumRealParameters}{0} & \Rval{3}{FullSignalNMSEdB}{3} dB & \Cval{3}{FLOPsPerSample}{0} & 14 módulos/raíces\\
PNNN N12 sparse & \Rval{8}{NumRealParameters}{0} & \Rval{8}{FullSignalNMSEdB}{3} dB & \Cval{8}{FLOPsPerSample}{0} & 14 módulos/raíces, 2 divisiones, 12 sigmoidales (hasta 12 exp)\\
\bottomrule
\end{tabularx}
\end{table}

El GMP usa operaciones complejas y envolventes compartibles. PN-IQ cuenta la
generación de features, productos por coeficiente, acumulaciones, normalización
y restauración. Coupled e independent full coinciden en 1062 FLOPs aunque tengan
200 y 358 parámetros: el acoplado necesita la aritmética cruzada equivalente a
productos complejos, mientras el independiente explota componentes Q nulas y
ajusta dos salidas. Parámetros y operaciones no son la misma magnitud.

La cifra principal de N12 sparse, 876 FLOPs/muestra, supone que la implementación
omite multiplicaciones por sus pesos exactamente nulos. Ejecutar las matrices
originales completas requiere 2252 FLOPs/muestra. Las redes añaden evaluaciones sigmoidales: H4 tiene
4 y N12, 12 por muestra. Ninguna de estas cifras demuestra automáticamente
tiempo MATLAB, coste FPGA, latencia o energía.

\begin{figure}[ht]
\centering
\begin{tikzpicture}
\begin{axis}[
 width=.78\linewidth,height=5.1cm,
 xlabel={FLOPs/muestra},ylabel={Full-signal NMSE (dB)},
 grid=both,
 tick label style={font=\small},label style={font=\small},
 xmin=750,xmax=2050,ymin=-39,ymax=-35.3]
\addplot[only marks,mark=star,mark size=4pt,color=teal]
 coordinates {(\PNFlops,\PNFull)};
\addplot[only marks,mark=square*,mark size=3pt,color=orange]
 coordinates {(\GMPMFlops,\GMPMFull)};
\addplot[only marks,mark=diamond*,mark size=3pt,color=redsoft]
 coordinates {(\SparseFlops,\SparseFull)};
\node[anchor=south west,font=\scriptsize] at (axis cs:\PNFlops,\PNFull) {PN-IQ (358)};
\node[anchor=north east,font=\scriptsize] at (axis cs:\GMPMFlops,\GMPMFull) {GMP matched (358)};
\node[anchor=south west,font=\scriptsize] at (axis cs:\SparseFlops,\SparseFull) {N12 sparse (358)};
\end{axis}
\end{tikzpicture}
\caption{Comparación principal. El conteo sparse supone que se omiten pesos nulos; las etiquetas indican parámetros reales.}
\label{fig:nmse-flops}
\end{figure}

\begin{figure}[ht]
\centering
\begin{tikzpicture}
\begin{axis}[
 ybar stacked,bar width=17pt,width=.72\linewidth,height=5cm,
 symbolic x coords={Omite ceros,Matrices completas},xtick=data,
 ylabel={FLOPs/muestra},ymajorgrids=true,
 tick label style={font=\small},label style={font=\small}]
\addplot[fill=teal!70] coordinates {(Omite ceros,876) (Matrices completas,2252)};
\end{axis}
\end{tikzpicture}
\caption{N12 sparse: 876 FLOPs/muestra omitiendo pesos nulos frente a 2252 si se ejecutan las matrices completas. Las operaciones adicionales se informan aparte.}
\label{fig:sparse-breakdown}
\end{figure}

\section{Resultados finales}

\begin{table}[ht]
\centering\scriptsize
\caption{Resultados numéricos de \texttt{comparison\_results.csv}.}
\begin{tabular}{@{}lrrrrrr@{}}
\toprule
Modelo & Entr. int. & Val. int. & Identificación & Señal completa & $\lambda$ & Soporte\\
\midrule
Complex GMP DOMP-100 & \Rval{0}{InternalTrainNMSEdB}{6} & \Rval{0}{InternalValidationNMSEdB}{6} & \Rval{0}{IdentificationNMSEdB}{6} & \Rval{0}{FullSignalNMSEdB}{6} & 0 & 100\\
Coupled PN-IQ & \Rval{1}{InternalTrainNMSEdB}{6} & \Rval{1}{InternalValidationNMSEdB}{6} & \Rval{1}{IdentificationNMSEdB}{6} & \Rval{1}{FullSignalNMSEdB}{6} & 0 & 100\\
Independent PN-IQ full & \Rval{2}{InternalTrainNMSEdB}{6} & \Rval{2}{InternalValidationNMSEdB}{6} & \Rval{2}{IdentificationNMSEdB}{6} & \Rval{2}{FullSignalNMSEdB}{6} & 0 & 100\\
Complex GMP matched & \Rval{3}{InternalTrainNMSEdB}{6} & \Rval{3}{InternalValidationNMSEdB}{6} & \Rval{3}{IdentificationNMSEdB}{6} & \Rval{3}{FullSignalNMSEdB}{6} & 0 & 179\\
Independent I/Q no PN & \Rval{4}{InternalTrainNMSEdB}{6} & \Rval{4}{InternalValidationNMSEdB}{6} & \Rval{4}{IdentificationNMSEdB}{6} & \Rval{4}{FullSignalNMSEdB}{6} & 0 & 100\\
Reduced independent PN-IQ & \Rval{5}{InternalTrainNMSEdB}{6} & \Rval{5}{InternalValidationNMSEdB}{6} & \Rval{5}{IdentificationNMSEdB}{6} & \Rval{5}{FullSignalNMSEdB}{6} & 0 & 100\\
PNNN H4 dense & \Rval{6}{InternalTrainNMSEdB}{6} & \Rval{6}{InternalValidationNMSEdB}{6} & \Rval{6}{IdentificationNMSEdB}{6} & \Rval{6}{FullSignalNMSEdB}{6} & -- & --\\
PNNN N12 dense & \Rval{7}{InternalTrainNMSEdB}{6} & \Rval{7}{InternalValidationNMSEdB}{6} & \Rval{7}{IdentificationNMSEdB}{6} & \Rval{7}{FullSignalNMSEdB}{6} & -- & --\\
PNNN N12 sparse & \Rval{8}{InternalTrainNMSEdB}{6} & \Rval{8}{InternalValidationNMSEdB}{6} & \Rval{8}{IdentificationNMSEdB}{6} & \Rval{8}{FullSignalNMSEdB}{6} & -- & --\\
\bottomrule
\end{tabular}
\end{table}

Las diferencias sobre la señal completa, redondeadas a tres decimales, son:
\begin{itemize}
\item Independent PN-IQ mejora 0.363 dB al GMP parameter-matched.
\item Independent PN-IQ mejora 2.634 dB a PNNN N12 sparse.
\item PNNN N12 sparse mejora 2.267 dB a H4 dense.
\item PNNN N12 dense mejora 0.731 dB a N12 sparse.
\item Reduced PN-IQ (200 parámetros) mejora 0.080 dB al GMP DOMP-100.
\end{itemize}
En coste, N12 sparse requiere 876 FLOPs/muestra si se omiten pesos nulos y 2252
si se ejecutan las matrices completas, mientras PN-IQ full usa 1062. La red
sparse ahorra 186 FLOPs bajo la primera implementación, pero requiere 1190 más
bajo la segunda y mantiene 12 evaluaciones sigmoidales adicionales.

\begin{figure}[ht]
\centering
\begin{tikzpicture}
\begin{axis}[
 width=.78\linewidth,height=5.2cm,
 xlabel={Parámetros reales activos},ylabel={Full-signal NMSE (dB)},
 grid=both,xmin=150,xmax=1100,ymin=-39,ymax=-33,
 tick label style={font=\small},label style={font=\small}]
\addplot[only marks,mark=*,mark size=2pt,color=navy]
 coordinates {(200,-37.9778360655445) (200,-37.9778360655442)
 (358,-38.5336473098535) (358,-38.1702651011503)
 (394,-37.8104589405615) (200,-38.0578594656762)
 (350,-33.6332342790055) (1046,-36.6310229561146)
 (358,-35.8997812260251)};
\addplot[only marks,mark=star,mark size=4pt,color=teal] coordinates {(\PNParams,\PNFull)};
\addplot[only marks,mark=square*,mark size=3pt,color=orange] coordinates {(\GMPMParams,\GMPMFull)};
\addplot[only marks,mark=diamond*,mark size=3pt,color=redsoft] coordinates {(\SparseParams,\SparseFull)};
\node[anchor=south west,font=\scriptsize] at (axis cs:\PNParams,\PNFull) {Independent PN-IQ};
\node[anchor=north east,font=\scriptsize] at (axis cs:\GMPMParams,\GMPMFull) {GMP matched};
\node[anchor=south west,font=\scriptsize] at (axis cs:\SparseParams,\SparseFull) {N12 sparse};
\end{axis}
\end{tikzpicture}
\caption{NMSE sobre la señal completa frente a parámetros activos. Los puntos proceden del CSV congelado.}
\label{fig:nmse-params}
\end{figure}

\section{Interpretación científica}

\subsection*{Lo demostrado en esta ejecución}
\begin{itemize}
\item La representación coupled PN-IQ reproduce el GMP complejo con error relativo
$\Rval{1}{RelativePredictionErrorToComplex}{12}$.
\item Relajar el acoplamiento mejora el full-signal NMSE del coupled de
\Rval{1}{FullSignalNMSEdB}{3} a \Rval{2}{FullSignalNMSEdB}{3} dB.
\item Con 358 parámetros en ambos modelos, Independent PN-IQ supera 0.363 dB al GMP matched.
\item El control sin PN queda 0.723 dB por detrás de Independent PN-IQ.
\item N12 sparse supera a H4 dense con presupuesto parecido, y los FLOPs se han
contado bajo una convención explícita.
\end{itemize}

\subsection*{Lo que sugieren, sin probar universalidad}
Las coordenadas PN parecen útiles cuando se libera la transformación I/Q. El
modelo lineal independiente puede aprovechar bien el régimen de pocos datos. La
red ancha y podada puede conservar conexiones que H4 no representa. Son hipótesis
consistentes con esta captura, no leyes generales.

\subsection*{Lo que no se ha demostrado}
No se demuestra superioridad universal, hardware, consumo, latencia, novedad
absoluta, ni generalización a otros PAs, señales o capturas. Tampoco se ha hecho
una campaña amplia de seeds neuronales. Los FLOPs son analíticos y no equivalen
a tiempo. Como la señal completa contiene identificación, tampoco constituye un
test independiente de generalización. No se presentan ACPR ni EVM demodulada
porque no hay una evaluación validada dentro de este resultado.

\section{Novedad y relación con trabajos previos}

Phase normalization, separación I/Q, representación real de productos
complejos, DOMP y redes phase-normalized no son ideas nuevas por separado. Hasta
donde se ha podido comprobar, la aportación potencialmente novedosa es la
combinación concreta: PN aplicado explícitamente a filas de una base GMP,
derivación coupled exactamente equivalente, extensión independent, controles de
coordenadas y presupuesto, comparación con una PNNN parameter-matched y análisis
de FLOPs. En una publicación conviene usar ``aportación potencialmente novedosa''
o \emph{to the best of our knowledge}, nunca ``first ever'' sin una revisión
bibliográfica sistemática.

\subsection{Decisiones de implementación y trazabilidad}

\begin{longtable}{@{}p{.19\linewidth}p{.21\linewidth}p{.25\linewidth}p{.27\linewidth}@{}}
\caption{Qué se hizo, motivo, alternativa y limitación.}\label{tab:decisions}\\
\toprule
Decisión & Implementación & Motivo/problema evitado & Alternativa y limitación\\
\midrule
\endfirsthead
\toprule Decisión & Implementación & Motivo/problema evitado & Alternativa y limitación\\\midrule
\endhead
Rotación por fila & \makecell[l]{\texttt{computePhaseNorm}\\\texttt{GMPRotation}} & Un marco común $\phi(n)$ para todos los taps & Rotar cada tap cambia el modelo; PN conserva dependencia del marco\\
$x(n)=0$ & $r=1$ & Evita división por cero & Otro valor unitario sería posible; convención fija\\
Extensión temporal & Índices periódicos compartidos & Igual dominio para todos & Recorte perdería bordes y cambiaría muestras\\
DC & Eliminación activada & Protocolo guardado común & Mantener DC sería otro experimento\\
Features PN & \makecell[l]{\texttt{buildPhaseNormalized}\\\texttt{IQRegressors}} & Reutiliza población y misma rotación & Generar señal global rotada usaría fases erróneas\\
Q nulas & \makecell[l]{\texttt{removeStructurally}\\\texttt{ZeroQFeatures}} & Elimina solo ceros exactos & Umbral aproximado introduciría reducción no controlada\\
Soporte coupled & DOMP-100 complejo compartido & Prueba equivalencia en mismo soporte & Reseleccionar podría confundir soporte y formulación\\
Parameter matched & DOMP complejo con 179 coeficientes & Iguala exactamente 358 parámetros reales & Igualar regresores ignoraría 2 parámetros/coeficiente\\
$\lambda$ & Split interno, rejilla congelada & Señal completa no selecciona & CV más amplia cuesta identificación; aquí gana 0\\
Escalado & Normas de columnas repetidas por pareja & Ridge equivalente y condición controlada & Escalas I/Q independientes rompen equivalencia\\
Restauración & $\widehat y=r^*\widehat y_{PN}$ & Devuelve el marco original & Omitirla evaluaría otra señal\\
PN reducido & \makecell[l]{\texttt{selectShared}\\\texttt{IQFeatures}} & 100 features comunes a dos salidas & Selección separada cambia presupuesto y soporte\\
Refit PNNN & Reinicio y épocas fijas con identificación completa & Todos los modelos ajustan las mismas 19,661 muestras & Usar solo entrenamiento interno dejaría menos datos para la red\\
Objetivo PNNN & 358 parámetros dinámicos & Igual a PN-IQ full & Un objetivo fijo puede quedar desparejado\\
Pruning & Global, no estructurado, por magnitud & Preserva conexiones más útiles globalmente & Structured sparsity facilitaría hardware, no se ensayó\\
Biases & Protegidos & Mantiene offsets y cuenta 14 & Podarlos cambia arquitectura efectiva\\
Freeze + fine-tune & Máscara fija & Evita que reaparezcan pesos & Regrowth sería otro algoritmo\\
FLOPs & Una cifra común y operaciones adicionales separadas & Mismo alcance de inferencia para todos & Medición temporal depende de plataforma\\
Predicciones & Señal completa en artefacto final & Mismo dominio para todos los modelos & Contiene identificación; no sustituye hold-out ni multi-captura\\
Asserts/tests & Invariantes de split, soporte y dimensiones & Detecta fugas y desajustes & No reemplaza verificación experimental externa\\
\bottomrule
\end{longtable}

\begin{table}[H]
\centering\small
\caption{Trazabilidad científica.}
\begin{tabularx}{\linewidth}{@{}lXX@{}}
\toprule
Concepto & Función/script principal & Resultado generado\\
\midrule
Phase normalization & \texttt{computePhaseNormGMPRotation.m}, \texttt{buildPhaseNormalizedIQRegressors.m} & Regresores PN y predicción restaurada\\
DOMP & \texttt{selectDOMPSupport.m}, \texttt{runPNGMPDOMPStudy.m} & \texttt{domp\_supports.mat}, history CSV\\
Coupled real & \texttt{complexToCoupledReal.m}, \texttt{fitPhaseNormGMPReal.m} & GMP/PN-IQ equivalente\\
Independent PN-IQ & \texttt{fitIndependentIQGMP.m}, \texttt{predictIndependentIQGMP.m} & PN-IQ full y reducido\\
Parameter matching & \texttt{runPNGMPDOMPStudy.m} & GMP matched congelado\\
Sparse PNNN & \texttt{refitFairPNNNDense.m}, \texttt{refitFairPNNNSparse.m} & modelos finales ajustados con identificación completa\\
NMSE & \texttt{nmseComplexDb.m} & columnas NMSE CSV\\
FLOPs & \texttt{countModelFLOPs.m}, \texttt{countPNNNFLOPs.m}, \texttt{countSparsePNNNFLOPs.m} & \texttt{complexity\_flops.csv}\\
Reporte & \texttt{run\_fair\_PNNN\_vs\_PNGMP\_DOMP.m} & CSV, MAT, PNG, PDF y predicciones full-signal\\
Tests & \texttt{run\_domp\_unit\_test.m}, \texttt{run\_fair\_comparison\_smoke\_test.m} & asserts PASS/FAIL\\
\bottomrule
\end{tabularx}
\end{table}

\section{Limitaciones y trabajo futuro}

\textbf{Limitaciones:} una captura principal; un número limitado de seeds de
PNNN; régimen fijo de identificación al 4\%; FLOPs analíticos; sin síntesis
hardware ni cuantización; sin ACPR/EVM demodulada validada; coste de entrenamiento
no igualado; dependencia de arquitectura, hiperparámetros y población GMP. La
red sparse se entrena sobredimensionada, mientras los lineales resuelven problemas
convexos; igualar parámetros activos no iguala proceso de optimización.

\textbf{Trabajo futuro:} varias capturas, señales y PAs; múltiples seeds;
barridos de datos; punto fijo; kernels DSP/FPGA; sparsity estructurada; latencia y
energía medidas; ACPR/EVM válidos; selección conjunta de soporte y acoplamiento;
y estudio de generalización fuera de distribución.

\section{Preguntas previsibles de los tutores}

\begin{enumerate}[leftmargin=*,label=\textbf{\arabic*.}]
\item \textbf{¿El PN-IQ acoplado es un modelo nuevo?} No como clase funcional:
es una representación real exactamente equivalente al GMP sobre el mismo soporte.
\item \textbf{¿Por qué da el mismo NMSE?} La rotación es unitaria, la matriz por
bloques representa el mismo producto complejo y la fase se restaura.
\item \textbf{¿Qué significa que los parámetros sean reales?} Son $\alpha$ y
$\beta$, componentes reales de coeficientes complejos; no implica coeficientes
GMP restringidos al eje real.
\item \textbf{¿Por qué el independiente tiene más libertad?} Sustituye la matriz
rotacional de dos grados por una transformación I/Q general de cuatro grados.
\item \textbf{¿Dónde entra exactamente PN?} En $r(n)$, aplicado a toda fila de
$U$ y a $y$, antes del ajuste, y deshecho con $r^*(n)$ tras predecir.
\item \textbf{¿Podría separarse I/Q sin PN?} Sí; por eso existe el control
Independent I/Q without PN.
\item \textbf{¿Qué demuestra ese control?} En esta captura, elegir coordenadas PN
mejora 0.723 dB al modelo independiente sin PN; no demuestra universalidad.
\item \textbf{¿Por qué DOMP?} Para reducir una población correlacionada mediante
selección y doble ortogonalización reproducibles.
\item \textbf{¿Por qué no OMP?} El estudio final congeló DOMP como único selector
para evitar mezclar algoritmos y porque sus tests quedaron validados.
\item \textbf{¿DOMP entra en FLOPs de inferencia?} No. Solo identifica soporte;
inferencia usa las columnas seleccionadas.
\item \textbf{¿Por qué 100 coeficientes GMP son 200 parámetros?} Cada complejo
$h=\alpha+j\beta$ contiene dos escalares reales libres.
\item \textbf{¿Por qué un GMP parameter-matched?} Para comprobar si la mejora se
explica solo por más parámetros frente al baseline de 200.
\item \textbf{¿Por qué la PNNN empieza con 1046 y acaba con 358?} N12 se entrena
ancha y el pruning conserva 344 pesos más 14 biases.
\item \textbf{¿Es completamente justa la comparación?} Es mucho más controlada
en identificación, dominio evaluado, métrica y parámetros activos; aún difieren
optimización y capacidad inicial. La evaluación completa contiene identificación
y no es un hold-out.
\item \textbf{¿Por qué N12 sparse mejora a H4 dense?} Puede conservar un patrón
distribuido dentro de una arquitectura ancha; aquí mejora 2.267 dB, sin garantía.
\item \textbf{¿Por qué no aparece la PNNN histórica de 1200 parámetros?} No cumple
el mismo presupuesto ni el mismo split; serviría como referencia, no comparación
principal.
\item \textbf{¿Cómo se cuentan los FLOPs de la red sparse?} La cifra principal
omite pesos exactamente nulos; ejecutar las matrices completas eleva el coste de
876 a 2252 FLOPs/muestra.
\item \textbf{¿La red sparse será realmente más rápida?} No necesariamente: hace
falta kernel sparse, y latencia depende de hardware, memoria y activaciones sigmoidales.
\item \textbf{¿Contribución defendible?} La derivación PN-GMP coupled, su extensión
independent con controles justos, y la comparación de parámetros/FLOPs.
\item \textbf{¿Qué no debemos afirmar?} ``Primero'', superioridad universal o
hardware, menor energía/latencia, o robustez multi-PA no medida.
\end{enumerate}

\section{Guion oral de cinco minutos}

\textbf{Problema (40 s).} ``Queremos modelar una señal compleja con memoria bajo
pocos datos de identificación y comparar GMP y PNNN sin mezclar splits ni
presupuestos. Usamos el mismo 4\% para ajustar todos los modelos y evaluamos
todos sobre la misma señal completa, que incluye ese 4\%.''

\textbf{Idea (55 s).} ``Rotamos cada fila del GMP con la fase conjugada de la
muestra actual. La forma coupled convierte el producto complejo en una matriz
real I/Q con dos parámetros por coeficiente y es exactamente equivalente. Después
liberamos el acoplamiento: I y Q tienen coeficientes independientes en ese marco.''

\textbf{Controles (50 s).} ``DOMP selecciona soporte solo en identificación. El
GMP matched controla el número de parámetros; I/Q sin PN controla las coordenadas;
PN-IQ reducido controla el presupuesto de 200. Coupled reutiliza exactamente el
soporte GMP.''

\textbf{PNNN (45 s).} ``Comparamos H4 dense, N12 dense y N12 sparse. La sparse
parte de 1046 parámetros y, con pruning global, biases protegidos y fine-tuning,
termina con 358 activos, igual que PN-IQ.''

\textbf{Resultados (55 s).} ``Sobre la señal completa, Independent PN-IQ alcanza
-38.534 dB; GMP matched, -38.170; N12 sparse, -35.900. PN aporta 0.723 dB frente
al control sin PN. N12 sparse mejora 2.267 dB a H4, pero pierde 0.731 frente a
N12 dense. Los tres modelos principales tienen 358 parámetros.''

\textbf{Coste y cierre (55 s).} ``PN-IQ usa 1062 FLOPs/muestra. N12 sparse usa
876 si se omiten pesos cero o 2252 si se ejecutan matrices completas, más 12
evaluaciones sigmoidales. Por tanto, PN-IQ gana NMSE en esta captura y la red gana el conteo
aritmético bajo ejecución sparse. La evaluación completa incluye identificación;
faltan hold-out multi-captura, múltiples seeds y validación hardware.''

\textbf{Versión de 30 segundos.} ``Reformulamos un GMP en coordenadas
phase-normalized. La variante coupled es algebraicamente el mismo GMP; al liberar
I/Q obtenemos una clase más flexible. Con 358 parámetros, Independent PN-IQ logra
-38.534 dB y PNNN N12 sparse -35.900 dB sobre la señal completa. PN-IQ cuesta
1062 FLOPs; la red sparse, 876 si se omiten ceros y 2252 con matrices completas,
más 12 evaluaciones sigmoidales. Es un resultado de una captura cuya evaluación incluye
identificación, no una victoria universal.''

\section{Resumen de una página para estudiar}

\begin{keybox}
\textbf{Fórmulas.}
$r=x^*/|x|$ ($r=1$ si $x=0$);
$U_{PN}=rU$; $y_{PN}=ry$; $\widehat y=r^*\widehat y_{PN}$.
$U_R=[U_I,-U_Q;U_Q,U_I]$.
$\mathrm{NMSE}_{dB}=10\log_{10}(\|y-\widehat y\|^2/\|y\|^2)$.
$P(H)=H(D+3)+2$, con $D=84$.
\end{keybox}

\begin{tabularx}{\linewidth}{@{}p{.22\linewidth}X@{}}
\toprule
Concepto & Mensaje memorizable\\
\midrule
Coupled & Dos parámetros por regresor; estructura rotacional; GMP exacto; error relativo $1.320\times10^{-11}$.\\
Independent & Coeficientes I/Q libres; más general, no equivalente; puede mejorar y sobreajustar.\\
Papel de PN & Marco de fase de la muestra actual; importa al romper acoplamiento; no rotar cada tap con su fase.\\
DOMP & Selección en identificación, doble ortogonalización, no entra en inferencia.\\
Protocolo & Identificación 19,661 (4.00004\%); selección interna 16,711/2,950; señal completa 491,520 e incluye identificación.\\
Parámetros & GMP-100=200; PN-IQ=358; GMP matched=358; H4=350; N12 sparse=358 activos.\\
Full-signal NMSE & PN-IQ $-38.534$; GMP matched $-38.170$; N12 sparse $-35.900$; H4 $-33.633$ dB.\\
FLOPs/muestra & PN-IQ 1062; GMP matched 1899; H4 876; N12 sparse 876 omitiendo ceros/2252 con matrices completas.\\
Special ops & PN-IQ 26; N12 sparse 54, incluidas 12 sigmoidales/exp de peor caso.\\
Sí puedo decir & Equivalencia coupled; mejora observada al liberar I/Q; controles parameter-matched y sin PN; coste analítico explícito.\\
No debo decir & Primero absoluto; mejor hardware; menor energía/latencia; generaliza a cualquier PA; sparse ideal garantizado.\\
\bottomrule
\end{tabularx}

\vspace{5mm}
\textbf{Tres diferencias clave:}
Independent PN-IQ supera 0.363 dB al GMP matched y 2.634 dB a N12 sparse;
N12 sparse supera 2.267 dB a H4. Omitiendo pesos nulos, N12 requiere 186 FLOPs
menos que PN-IQ, pero con matrices completas requiere 1190 más.

\textbf{Respuesta final si preguntan ``¿quién gana?''}
``En NMSE de esta captura gana Independent PN-IQ. En FLOPs/muestra gana la PNNN
sparse si la implementación omite ceros; si ejecuta matrices completas, su coste
sube a 2252. Sin plataforma ni métrica prioritaria no hay ganador absoluto.''

{\small\raggedright
\subsection*{Fuentes congeladas}
Las tablas proceden de \path{comparison_results.csv},
\path{parameter_summary.csv}, \path{complexity_flops.csv} y
\path{selected_hyperparameters.csv}. La partición y configuración se verifican
en \path{split_indices.mat} y \path{comparison_config.mat}; soportes en
\path{domp_supports.mat}; predicciones finales en
\path{full_signal_predictions.mat}. Todos pertenecen a
\path{results/full_signal_domp_comparison/20260714_102201}.

\medskip
\textbf{Referencia DOMP.}
\begingroup
\renewcommand{\section}[2]{}
\begin{thebibliography}{1}
\bibitem{becerra_domp}
J. A. Becerra et al.,
\newblock ``A Doubly Orthogonal Matching Pursuit Algorithm for Sparse
Predistortion of Power Amplifiers.''
\newblock Referencia DOMP indicada en el proyecto; los artefactos locales no
incluyen metadatos bibliográficos completos.
\end{thebibliography}
\endgroup
}

\end{document}
